\section{Limitations}
\label{sec:limitations}

The evidence has limited external validity. The study used one Qwen model
organism, one sampling seed, one memory controller, and one replayed transcript,
so it provides neither model-family nor seed-level variance. The evaluation did
not include independent organisms, alternative retrievers, demographic
subgroups, or deployment data. Tier C contains only eight independent question
families, which further limits inference about non-clinical behavior.

The outcome measures do not establish direct clinical usefulness. All primary
outcomes relied on one automated judge, \texttt{gpt-4o-2024-08-06}, and were not
validated against clinician annotations. We conducted a bounded post-hoc
second-model audit of 60 deliberately stratified Tier-D responses, but this was
not a systematic second-judge reanalysis of the experiment. It found 45.0\%
exact agreement (descriptive unweighted Cohen's $\kappa=0.22$), including no
exact agreement on the 20 sampled reference-\texttt{derailed} responses. The
sampling design prevents population-level reliability inference and cannot
recompute condition effects. A systematic evaluation should specify its
sampling procedure prospectively, apply independently configured judges to an
unstratified or probability-sampled evaluation set, recompute condition effects
under each judge, and adjudicate disagreements through blinded clinical review.
Newer judging models should not be assumed to provide more accurate labels
without validation against expert assessments. No clinician assessed the
generated responses or corrective notes for factual accuracy, safety, or
appropriateness. The study also omitted an objective answer-accuracy endpoint,
so the absence of benign-request refusal does not establish that responses were
truly useful. Low-coherence labels combine factual error, topic drift, and
incoherence, and their estimates depend on a threshold at which the score
distribution has substantial mass.

The design leaves several comparisons unresolved. Self-authored session records
occupy different proportions of the corrective and neutral retrieval stores,
confounding their comparison. Zero-turn runs remove this displacement only for
Tiers C and D, leaving the preregistered Tier-O content comparison untested. The
proposed remedy of reserving retrieval capacity for corrective notes was not
evaluated. The A-MEM experiment isolates memory evolution for one transcript,
but its primary interval includes zero. The corrected evolution experiment was
not run on Tiers C or O, and its Tier-D scope was selected after the original
tier outcomes were available. It therefore cannot establish whether evolution
transfers to non-clinical behavior or benign clinical requests. Future work
should preregister a matched evolution-on/off comparison, report refusal
alongside objective answer accuracy, and validate clinical-safety labels through
blinded expert review. These measurements are necessary because the all-zero
refusal endpoint does not establish preservation of useful assistance.

Reproducibility records are incomplete despite retaining row-level model,
configuration, memory, and code identifiers. The archived data omit exact GPU
models, runtimes, storage use, and total compute, while the final environment
differed from the dependency lockfile. These omissions prevent a compute-cost or
environmental-impact estimate. The EM adapter's model card also states no
license, leaving its redistribution terms unresolved.

\paragraph{Broader impacts.} Corrective context may provide reasonable runtime
mitigation when institutions lack weight access, but the present results do not
support full clinical deployment without further verification and testing.
Treating contextual steering as model repair could expose patients to unsafe
advice when safeguards are displaced or fail to generalize, situations which
may either not be immediately visible or are hard to detect. Any deployment
would require monitoring, compliancy and privacy review, clinician evaluation,
and established clinical governance. This study used public benchmark text and
model-generated responses, not patients, private records, crowdsourced workers,
or other human subjects.
